\section*{Capítulo \ref{ch:g12:early-earth} --- A Terra primitiva e a sua atmosfera}

\begin{solution}{exo:g12:early-earth:1}
(a) a \omterm{def:g12:early-earth:primitive-atmosphere}{desgaseificação}; (b) um \omterm{def:g12:early-earth:stromatolite}{estromatólito}; (c) uma \omterm{def:g12:early-earth:bif}{formação ferrífera
bandada}; (d) as \omterm{def:g12:early-earth:red-beds}{camadas vermelhas}; (e) o \omterm{def:g12:early-earth:goe}{Grande Evento de Oxidação}.
\end{solution}

\begin{solution}{exo:g12:early-earth:2}
Os zircões que mostram água líquida (4.4 bilhões de anos); os primeiros
\omterm{def:g12:early-earth:stromatolite}{estromatólitos} da Pilbara (3.43); as \omterm{def:g12:early-earth:bif}{formações ferríferas bandadas} das
Hamersley Ranges (mais de 2.5); o fim do \omterm{def:g12:early-earth:goe}{Grande Evento de Oxidação} (2.06); o
fim da maioria das \omterm{def:g12:early-earth:bif}{formações ferríferas bandadas} (1.8).
\end{solution}

\begin{solution}{exo:g12:early-earth:3}
Em $t = 0$, $L = L_\odot/(1 + 0.4) = L_\odot/1.4 \approx 0.714\,L_\odot$:
cerca de \qty{29}{\percent} menos brilhante, quase um terço a menos que hoje.
\end{solution}

\begin{solution}{exo:g12:early-earth:4}
Que o ar continha então quase nenhum oxigênio. A pirita é destruída pelo
oxigênio; os seus grãos só podiam ser rolados pelos rios e depositados
intactos se o ar e a água da superfície não o tivessem.
\end{solution}

\begin{solution}{exo:g12:early-earth:5}
O vapor de água, que se condensou no oceano, e o gás carbônico, que ficou no
ar em quantidades muito maiores que hoje e aqueceu a Terra pelo seu efeito
estufa.
\end{solution}

\begin{solution}{exo:g12:early-earth:6}
Há 3 bilhões de anos: $1 - t/t_\odot = 3/4.567 \approx 0.657$, $L \approx
1/1.263 \approx 0.792\,L_\odot$, $T \approx 255 \times 0.792^{1/4} \approx
\qty{241}{K}$. Há 2.45 bilhões de anos: $1 - t/t_\odot \approx 0.536$, $L
\approx 1/1.215 \approx 0.823\,L_\odot$, $T \approx 255 \times 0.823^{1/4}
\approx \qty{243}{K}$.
\end{solution}

\begin{solution}{exo:g12:early-earth:7}
\qty{1}{kg} de ferro são $1000/55.85 \approx 17.9$ mols; uma molécula de
oxigênio capta quatro íons de ferro, então são $17.9/4 \approx 4.48$ mols de
\ce{O2}, ou seja, $4.48 \times 32.00 \approx \qty{143}{g}$ de oxigênio.
\end{solution}

\begin{solution}{exo:g12:early-earth:8}
(1) $2.46 - 2.06 = 0.40$ bilhão de anos, 400 milhões de anos. (2) De 4.54
a 2.46 bilhões de anos atrás: $2.08/4.54 \approx \qty{46}{\percent}$ da
história da Terra. (3) Cerca de $3.43 - 2.46 \approx 1$ bilhão de anos: o
oxigênio produzido pelos micróbios foi consumido por muito tempo antes de
poder ficar no ar.
\end{solution}

\begin{solution}{exo:g12:early-earth:9}
A \omterm{def:g12:early-earth:bif}{formação ferrífera bandada} mostra um mar profundo sem oxigênio, que podia
transportar ferro dissolvido, sob um mar raso e iluminado com oxigênio
bastante para precipitá-lo. A camada vermelha, depositada nos continentes,
mostra oxigênio no ar e na água da superfície, o bastante para enferrujar o
ferro dos seus grãos.
\end{solution}

\begin{solution}{exo:g12:early-earth:10}
Em Hamelin Pool, Shark Bay, na Austrália Ocidental, cuja água é duas vezes
mais salgada que a água do mar comum. Um tapete de micróbios, sobretudo
cianobactérias, retém e liga grãos de sedimento e faz crescer carbonato; um
novo tapete cresce por cima, e assim por diante, camada após camada.
\end{solution}

\begin{solution}{exo:g12:early-earth:11}
(1) $21 \times 10^{-5} = \qty{0.00021}{\percent}$. (2) $21/10 =
\qty{2.1}{\percent}$.
\end{solution}

\begin{solution}{exo:g12:early-earth:12}
A \omterm{def:g10:habitable-planet:equilibrium-temperature}{temperatura de equilíbrio} não muda se $L(1 - A')$ for igual a $1 - A$:
$1 - A' = 0.706/0.741 \approx 0.953$, logo $A' \approx 0.047$. A Terra teria
precisado refletir menos de \qty{5}{\percent} da luz do Sol, quase como um
corpo negro, o que um planeta oceânico com nuvens e um pouco de gelo não
consegue. A outra saída é um efeito estufa mais forte.
\end{solution}

\begin{solution}{exo:g12:early-earth:13}
(1) $273 - 237 = \qty{36}{K}$ para chegar ao congelamento; $288 - 237 = \qty{51}{K}$ para
chegar à média de hoje. (2) O efeito estufa de hoje acrescenta $288 - 255 =
\qty{33}{K}$: a Terra jovem precisava de um mais forte. (3) \qty{0.1}{bar} é um
décimo da pressão da atmosfera inteira de hoje, só para o gás carbônico.
\end{solution}

\begin{solution}{exo:g12:early-earth:14}
(1) Antes de 2450 milhões de anos atrás a assinatura existe, então o ar
continha menos de $10^{-5}$ do oxigênio de hoje; depois de 2090 milhões de
anos ela está ausente, então o ar continha mais. (2) Sim: o \omterm{def:g12:early-earth:goe}{Grande Evento de
Oxidação} vai de cerca de 2460 a 2060 milhões de anos atrás, o mesmo intervalo.
\end{solution}

\begin{solution}{exo:g12:early-earth:15}
(1) Cerca de $4567 - 60 = 4507$ milhões de anos atrás. (2) No máximo $4507 - 4404 =
103$ milhões de anos depois do impacto: a Terra tinha uma crosta e água
menos de cerca de cem milhões de anos depois do impacto.
\end{solution}

\begin{solution}{pb:g12:early-earth:1}
\textbf{1.} $4.567 - 4.0 = 0.567$ bilhão de anos.

\textbf{2.} $1 - 0.567/4.567 \approx 1 - 0.124 = 0.876$.

\textbf{3.} $L = 1/(1 + 0.4 \times 0.876) = 1/1.350 \approx 0.741$ da
de hoje.

\textbf{4.} $0.741 \times 1361 \approx \qty{1008}{W/m^2}$.

\textbf{5.} $1 - 0.741 \approx 0.26$: \qty{26}{\percent} menos brilhante.

\textbf{6.} $T_{\mathrm{eq}} = \left(S(1-A)/(4\sigma)\right)^{1/4}$.

\textbf{7.} Com $A$ e $\sigma$ fixos, $T_{\mathrm{eq}}$ é uma constante vezes
$S^{1/4}$, logo $T'/T = (S'/S)^{1/4}$.

\textbf{8.} $0.741^{1/4} = \sqrt{\sqrt{0.741}} \approx \sqrt{0.861} \approx
0.928$.

\textbf{9.} $255 \times 0.928 \approx \qty{237}{K}$.

\textbf{10.} $237 - 273 \approx \qty{-36}{\celsius}$.

\textbf{11.} Que havia água líquida na superfície, um oceano, há 4.4
bilhões de anos.

\textbf{12.} Em águas rasas, como os \omterm{def:g12:early-earth:stromatolite}{estromatólitos} vivos de hoje.

\textbf{13.} As \omterm{def:g12:early-earth:bif}{formações ferríferas bandadas} das Hamersley Ranges,
depositadas num fundo marinho.

\textbf{14.} Sim: \qty{237}{K} fica muito abaixo de \qty{273}{K}; com a
atmosfera de hoje, a superfície média teria ficado congelada.

\textbf{15.} As rochas mostram água líquida sob um Sol fraco demais para
mantê-la líquida com a atmosfera de hoje: o \omterm{def:g12:early-earth:faint-young-sun}{paradoxo do Sol jovem fraco}.

\textbf{16.} $288 - 255 = \qty{33}{K}$.

\textbf{17.} Pelo menos $273 - 237 = \qty{36}{K}$, mais que hoje.

\textbf{18.} $A' = 1 - 0.706/0.741 \approx 0.047$: um planeta quase negro,
não plausível.

\textbf{19.} O gás carbônico, muito mais abundante que hoje; o metano,
produzido pelos primeiros micróbios, pode ter ajudado.

\textbf{20.} A \omterm{def:g10:habitable-planet:equilibrium-temperature}{temperatura de equilíbrio} da Terra há quatro bilhões de anos
era de \textbf{cerca de 237 K}, uns \qty{36}{\celsius} abaixo do ponto de
congelamento.
\end{solution}
